\section{Introduction}

The method of least squares \index{least squares} is one of Gauss's most influential contributions to science. He developed it for analyzing astronomical observations (especially the orbit of Ceres) and published it in \textit{Theoria Motus Corporum Coelestium} (1809). The method minimizes the sum of squared residuals and provides the best linear unbiased estimator under the Gauss-Markov theorem \index{Gauss-Markov theorem}.

In this chapter, we cover:
\begin{itemize}
    \item Linear least squares: $Ax \approx b$
    \item Normal equations: $A^T A x = A^T b$
    \item QR decomposition solution
    \item Weighted least squares
    \item Regularization (Tikhonov, ridge regression \index{ridge regression})
    \item Nonlinear least squares: Gauss-Newton, Levenberg-Marquardt \index{Levenberg-Marquardt algorithm}
    \item Gauss's application to the orbit of Ceres
\end{itemize}

\section{Linear Least Squares}

Given $A \in \R^{m \times n}$ ($m \ge n$) and $b \in \R^m$, find $x$ minimizing $\|Ax - b\|_2^2$.

\begin{theorem}[Normal Equations \index{normal equations}]
The minimizer satisfies $A^T A x = A^T b$. If $A$ has full column rank, $A^T A$ is invertible and $x = (A^T A)^{-1} A^T b$.
\end{theorem}

\begin{theorem}[Gauss-Markov Theorem \index{Gauss-Markov theorem}]
The least squares estimator is the best linear unbiased estimator (BLUE): among all linear estimators $\hat{x} = Cy$ with $E[\hat{x}] = x$, the LS estimator has minimum covariance.
\end{theorem}

\begin{definition}[Residuals \index{residuals}]
The residual vector $r = b - Ax$ measures the discrepancy between observed and predicted values.
\end{definition}

\section{QR Solution}

Using $A = QR$ with $Q^T Q = I$, $R$ upper triangular:
\[
\|Ax - b\|_2 = \|QRx - b\|_2 = \|Rx - Q^T b\|_2.
\]
Solve $Rx = Q^T b$ by back substitution. This is more stable than normal equations.

\begin{definition}[Condition Number \index{condition number}]
The condition number $\kappa(A) = \|A\| \cdot \|A^{-1}\|$ measures sensitivity to perturbations. For least squares, $\kappa(A^TA) = \kappa(A)^2$.
\end{definition}

\section{Weighted Least Squares}

Minimize $\|W^{1/2}(Ax - b)\|_2^2$ for symmetric positive definite weight matrix $W$. Solution: $(A^T W A)x = A^T W b$.

\begin{definition}[Information Matrix \index{information matrix}]
The matrix $A^T W A$ is the information matrix, measuring the amount of information about $x$ contained in the observations.
\end{definition}

\section{Tikhonov Regularization}

For ill-conditioned problems, minimize $\|Ax - b\|_2^2 + \lambda \|x\|_2^2$:
\[
(A^T A + \lambda I)x = A^T b.
\]
The regularization parameter $\lambda$ controls the bias-variance tradeoff.

\begin{definition}[Ridge Regression \index{ridge regression}]
Tikhonov regularization with $\lambda > 0$ is equivalent to ridge regression in statistics, adding a penalty on the $L_2$ norm of coefficients.
\end{definition}

\section{Nonlinear Least Squares}

Minimize $\|f(x)\|_2^2$ where $f: \R^n \to \R^m$. The Gauss-Newton method \index{Gauss-Newton method} linearizes:
\[
x_{k+1} = x_k - (J_k^T J_k)^{-1} J_k^T f(x_k)
\]
where $J_k$ is the Jacobian of $f$ at $x_k$.

The Levenberg-Marquardt algorithm \index{Levenberg-Marquardt algorithm} adds damping:
\[
x_{k+1} = x_k - (J_k^T J_k + \lambda I)^{-1} J_k^T f(x_k).
\]
This combines gradient descent ($\lambda$ large) with Gauss-Newton ($\lambda$ small).

\section{Gauss and the Orbit of Ceres}

In 1801, asteroid Ceres was discovered but lost. Gauss used least squares on three observations to predict its orbit, enabling its rediscovery. He solved a nonlinear least squares problem for the six orbital elements.

\begin{definition}[Orbit Determination \index{orbit determination}]
The problem of finding orbital elements from observations is called orbit determination. Gauss's method is a preliminary orbit determination technique.
\end{definition}

\section{Python Implementation}

In \texttt{python/least\_squares.py}:

\begin{lstlisting}[style=python, caption={Key functions in python/least\_squares.py}]
least_squares_normal(A, b)      # Normal equations
least_squares_qr(A, b)          # QR decomposition
weighted_least_squares(A, b, W) # Weighted LS
ridge_regression(A, b, lam)     # Tikhonov regularization \index{Tikhonov regularization}
gauss_newton(f, jac, x0)        # Gauss-Newton
levenberg_marquardt(f, jac, x0) # Levenberg-Marquardt
ceres_orbit(data)               # Gauss's Ceres problem
\end{lstlisting}

\section{Numerical Examples}

\begin{example}[Linear Least Squares]
\begin{lstlisting}[style=python]
>>> from least_squares import least_squares_qr
>>> import numpy as np
>>> A = np.array([[1, 1], [1, 2], [1, 3], [1, 4]], dtype=float)
>>> b = np.array([2.1, 3.9, 6.2, 7.8])
>>> x = least_squares_qr(A, b)
>>> print(x)  # [0.85, 1.75] -> y = 0.85 + 1.75x
\end{lstlisting}
\end{example}

\section{Exercises}

\begin{exercise}
Derive the weighted least squares \index{weighted least squares} solution from first principles.
\end{exercise}

\begin{exercise}
Implement the QR updating/downdating algorithms for recursive least squares.
\end{exercise}

\begin{exercise}
Apply Levenberg-Marquardt to fit an exponential model $y = a e^{bx} + c$.
\end{exercise}

\begin{exercise}
Reconstruct Gauss's Ceres orbit determination \index{orbit determination} using modern observations.
\end{exercise}

\begin{exercise}
Implement the iterative refinement algorithm for least squares solutions.
\end{exercise}



\section{Exercise Solutions}

\begin{solution}
The normal equations $A^TAx = A^Tb$ arise from minimizing $\|Ax - b\|^2$. The Hessian $2A^TA$ is positive definite if $A$ has full column rank, making this a convex optimization problem. The solution is $x = (A^TA)^{-1}A^Tb$ when $A^TA$ is invertible.
\end{solution}

\begin{solution}
Using the QR decomposition $A = QR$ with $Q$ orthogonal, the least squares problem becomes $\min_x \|Rx - Q^Tb\|^2$, which is solved by back substitution. This is numerically more stable than forming $A^TA$ directly, as it avoids squaring the condition number.
\end{solution}

\begin{solution}
The residual $r = b - Ax$ is orthogonal to the column space of $A$: $A^Tr = 0$. The covariance of the least squares estimator is $\operatorname{Cov}(\hat{x}) = \sigma^2(A^TA)^{-1}$ assuming $\operatorname{Cov}(b) = \sigma^2 I$. Gauss showed that if errors are normally distributed, least squares gives the maximum likelihood estimator.
\end{solution}

\begin{solution}
Iterative refinement solves $A\delta = r$ where $r = b - A\hat{x}$ is the residual, then updates $\hat{x} \leftarrow \hat{x} + \delta$. Each iteration reduces the error by a factor of $\kappa(A)\epsilon$. In practice, 2--3 iterations suffice to reach machine precision for well-conditioned systems.
\end{solution}

\begin{solution}
The QR updating algorithm maintains the QR decomposition as new data arrives. For rank-1 updates, the Givens rotation structure allows $O(n^2)$ updates rather than $O(n^3)$ recomputation. Downdating removes a row using hyperbolic rotations.
\end{solution}


\section{References}

\begin{itemize}
    \item \cite{gauss-theoria} -- Gauss, \textit{Theoria Motus Corporum Coelestium} (1809)
    \item \cite{golub-vanloan} -- Ch. 5 (Least Squares)
    \item \cite{bjorck} -- Björck, \textit{Numerical Methods for Least Squares Problems}
    \item \cite{nocedal-wright} -- Nocedal \& Wright, \textit{Numerical Optimization}
\end{itemize}